% AUTO-GENERATED durch diagram_generator (E-2c, Pareto-/Tradeoff-Streuung p50 vs p99; lang=en)
% Ein Punkt je (Konfiguration x Op-Art): x = Median-Latenz, y = Tail-Latenz (p99). Beide Groessen
% stehen SO im WIDE-Schema -- nichts ist erfunden, nichts aggregiert. Diagonale y=x = kein
% Tail-Aufschlag; je weiter ein Punkt darueber liegt, desto teurer sein Ausreisser-Verhalten.
% Achsen LOG.
\begin{figure}[!htbp]
\centering
\resizebox{\textwidth}{!}{%
\begin{tikzpicture}
\definecolor{tradeoff0}{RGB}{217,76,76}
\definecolor{tradeoffdiag}{RGB}{120,120,120}
\begin{axis}[
    width=0.7800\textwidth,
    height=0.4000\textheight,
    scale only axis,
    title={Pareto scatter: median vs tail latency},
    xlabel={median latency p50 (ns/op)},
    ylabel={tail latency p99 (ns/op)},
    grid=both,
    enlargelimits=0.05,
    xmode=log,
    ymode=log,
    legend style={at={(1.03,1)},anchor=north west,font=\tiny,legend cell align=left},
    mark size=1.6pt,
]
% E-2c-DIAGONALE y=x: Referenz "kein Tail-Aufschlag". Punkte darueber zahlen Aufschlag.
\addplot[sharp plot,no marks,dashed,tradeoffdiag,forget plot] coordinates {(750.0000,750.0000) (1170.0000,1170.0000)};
\addplot[only marks,mark=*,color=tradeoff0] coordinates {(750.0000,1170.0000)};
\addlegendentry{k\_ary}
\end{axis}
\end{tikzpicture}
}%
\caption{Pareto scatter: median vs tail latency (one point per configuration and operation kind; diagonal = no tail surcharge)}
\end{figure}
